\documentclass[11pt]{article}
\usepackage[utf8]{inputenc}
\usepackage[T1]{fontenc}
\usepackage{amsmath,amssymb,amsthm}
\usepackage[margin=1.1in]{geometry}
\usepackage{hyperref}
\usepackage{booktabs}

\newtheorem{theorem}{Theorem}
\newtheorem{proposition}[theorem]{Proposition}
\newtheorem{lemma}[theorem]{Lemma}
\theoremstyle{definition}
\newtheorem{remark}[theorem]{Remark}

\title{Searching the recursion tree improves the lower bound\ on the Shannon capacity of $C_7$}
\author{Oleksii Stavriianov\ \small ORCID 0009-0002-8504-3626}
\date{}

\begin{document}
\maketitle

\begin{abstract}
The recent sequence of improvements to the lower bound on the Shannon capacity of the seven-cycle
shares a common shape: a library of admissible combining rules is applied along a fixed recursion
tree. The rules have been studied carefully; the tree has not. We show that optimising the tree,
using no data beyond what is already published, gives
\[
  \Theta(C_7) \;\ge\; 3.258834362237710794\ldots,
\]
improving the bound $3.258832620353266309\ldots$ of Tandon. The construction is
$K_{4a}$ applied to four copies of $S_{3b}(n_6, n_{11}, n_{11})$, where $n_6$ and $n_{11}$ are
Tandon's own published cardinality vectors and $S_{3b}$, $K_{4a}$ are rules of Buys, Polak and
Zuiddam: the only change is that the third input is a second copy of $n_{11}$ rather than $n_8$.
Among all $16\,384$ combinations of the three published blocks, the published choice is third
by value.

We also record what does \emph{not} help, since several natural avenues can be closed
exhaustively rather than merely explored. The base gadget of Polak and Schrijver admits exactly
eight private pairs and this is a structural ceiling; its neutral auxiliary part cannot exceed
$322$ over the entire automorphism orbit; nine private pairs do exist, on other $367$-sets, but
are incompatible with a full-sized auxiliary set, with a quantitative obstruction; and the
published twist is optimal over all $64\,538\,880$ automorphisms. Finally we note two traps that
internal consistency does not catch.
\end{abstract}

\section{Introduction}

Shannon's zero-error capacity $\Theta(G)$~\cite{Shannon} is known for no odd cycle beyond $C_5$, where
Lov\'asz's bound~\cite{Lovasz} settles it. For $C_7$ the gap remains wide: the best upper bound is still
$\vartheta(C_7) = 3.3176672\ldots$, while the lower bound has been improved four times in the
space of two months, each time in the seventh significant digit.

The improvements share a structure. One fixes a base object in $C_7^{\boxtimes 5}$, a set of
admissible rules for combining such objects, and a recursion tree describing which rule is applied
to which intermediates. Itty, Rosin, Carstensen and Reichman found a $134\,753$-element independent set in
$C_7^{\boxtimes 10}$~\cite{Itty}; Gao recast it as a binary product lemma over
gadgets built from the $367$-set of Polak and Schrijver~\cite{PS} and reached
$3.2587891539\ldots$~\cite{Gao}; Buys, Polak and Zuiddam (BPZ)~\cite{BPZ1,BPZ2} replaced the six numbers of a Gao gadget by
seven families with prescribed separations, verified a library of ten combining rules and three
terminal codes in Lean, and reached $3.2588279859\ldots$; Tandon added a heterogeneous refinement
of the binary product and reached $3.2588326203\ldots$~\cite{Tandon}.

In every one of them the recursion tree is fixed by hand. Tandon is explicit: the ordered split
$25 = 6 + 11 + 8$, the rule $S_{3b}$ and the terminal code $K_{4a}$ are kept unchanged from BPZ,
and only the representations sitting at the nodes are replaced. The rules have been scrutinised;
the tree has not been searched.

This note searches it. The result is a better bound obtained from published data alone, and --
more usefully -- evidence that the tree is a systematically underexploited parameter rather than a
lucky one.


\paragraph{Scope.}
The improvement is in the seventh significant digit, as were the four before it, and the distance
to the upper bound is unaffected: $\vartheta(C_7) - \Theta(C_7)$ is still bounded only by
$2.9 \times 10^{-2}$, four orders of magnitude larger than everything this line of work has gained
in total. Nor is the construction deep. It is one substitution in someone else's recursion, and
stating it here makes it reproducible in an afternoon by anyone holding the same data; we expect
it to be superseded quickly, and by the authors whose data it uses. What we hope survives is the
observation that the tree is a free parameter, which applies to whatever construction comes next,
and the exhaustive verifications of Section~5, which close doors rather than open them but were
expensive to establish and are cheap to reuse.

\section{The BPZ framework, briefly}

We use the seven-family formulation of BPZ~\cite{BPZ2} as presented by Tandon~\cite{Tandon}. Let
$L = \{B, N, A, D, O, H, V\}$ carry the symmetric separation relation $\perp$ given by
$B \perp \{O,H,V\}$, $N \perp \{A,D,O,H,V\}$, $A \perp \{N,D,H\}$, $D \perp \{N,A,V\}$,
$O \perp \{B,N\}$, $H \perp \{B,N,A\}$, $V \perp \{B,N,D\}$. A \emph{seven-family
representation} in a graph $G$ is a collection $F = \{F_\lambda\}_{\lambda \in L}$ of independent
sets with $F_\lambda$ separated from $F_\mu$ whenever $\lambda \perp \mu$, and its
\emph{cardinality vector} is $(|F_B|, \ldots, |F_V|)$ in the order $(B,N,A,D,O,H,V)$.

An $m$-ary \emph{combining rule} is a family $T = \{T_\lambda\}$ with $T_\lambda \subseteq L^m$,
admissible when any two words of one component are separated in some coordinate and words of
components with separated labels are too. For admissible $T$ and representations
$F^{(1)}, \ldots, F^{(m)}$ the sets
$F^{\mathrm{out}}_\lambda = \bigcup_{(\lambda_1,\ldots,\lambda_m) \in T_\lambda}
 F^{(1)}_{\lambda_1} \times \cdots \times F^{(m)}_{\lambda_m}$
form a representation in the strong product, with
\begin{equation}\label{eq:card}
  |F^{\mathrm{out}}_\lambda| \;=\; \sum_{(\lambda_1,\ldots,\lambda_m) \in T_\lambda}
  \prod_{i=1}^m |F^{(i)}_{\lambda_i}|.
\end{equation}
A \emph{terminal code} is a set of $m$-tuples whose Cartesian products are pairwise separated; its
union is a single independent set, of size given by the same formula. The BPZ library holds ten
combining rules $S_{2a}, S_{2b}, S_{3a}, \ldots, S_{3h}$ and three terminal codes
$K_{3a}, K_{4a}, K_{4b}$, all verified admissible in Lean. Gao's binary product is the rule
$S_{2a}$.

Nothing in \eqref{eq:card} requires the inputs to be distinct. This is not a loophole: the
published constructions already use it, applying $K_{4a}$ to four copies of one representation.

\section{The improvement}

Tandon's Appendix B.4~\cite{Tandon} prints three cardinality vectors, of representations in
$C_7^{\boxtimes 30}$, $C_7^{\boxtimes 55}$ and $C_7^{\boxtimes 40}$ respectively:
\[
  n_6, \qquad n_{11}, \qquad n_8,
\]
and forms $n_{25} = S_{3b}(n_6, n_{11}, n_8)$, a representation in $C_7^{\boxtimes 125}$, then
applies $K_{4a}$ to four copies to obtain an independent set in $C_7^{\boxtimes 500}$ of $257$
digits, giving $3.258832620353266309\ldots$.

\begin{theorem}\label{thm:main}
Let $v = S_{3b}(n_6, n_{11}, n_{11})$ and let $M = K_{4a}(v,v,v,v)$. Then $M$ is the cardinality
of an independent set in $C_7^{\boxtimes 560}$, it has $288$ digits, and
\[
  \Theta(C_7) \;\ge\; M^{1/560} \;=\; 3.258834362237710794\ldots
\]
\end{theorem}

\begin{proof}
$n_6$ and $n_{11}$ are cardinality vectors of seven-family representations in
$C_7^{\boxtimes 30}$ and $C_7^{\boxtimes 55}$; $S_{3b}$ is an admissible ternary combining rule
and $K_{4a}$ an admissible terminal code, both from the Lean-verified BPZ library. By
\eqref{eq:card}, $v$ is the cardinality vector of a representation in
$C_7^{\boxtimes 30+55+55} = C_7^{\boxtimes 140}$, and $K_{4a}$ applied to four copies of it gives
an independent set in $C_7^{\boxtimes 560}$. The arithmetic is elementary and exact; the value is
recorded in Appendix~A.
\end{proof}

The only difference from the published construction is the third input. The split
$28 = 6 + 11 + 11$ replaces $25 = 6 + 11 + 8$, and the dimension grows from $500$ to $560$.

\section{The tree is a systematically underexploited parameter}

One improvement could be luck. Three observations suggest otherwise.

\paragraph{The published choice is not the best among its own inputs.}
Drawing three blocks with repetition from $\{n_6, n_8, n_{11}\}$ and the base, in both
orientations, under all eight ternary rules and all three terminal codes, and also reading each
product directly as an independent set of size $|F_B| + |F_O|$, gives $16\,384$ candidates. The
ranking begins

\begin{center}
\begin{tabular}{llr}
\toprule
construction & dimension & bound \\
\midrule
$K_{4a}(S_{3b}(n_6, n_{11}, n_{11}))$ & $560$ & $3.2588343622377$ \\
$K_{4a}(S_{3b}(n_8, n_{11}, n_{11}))$ & $600$ & $3.2588331674194$ \\
$K_{4a}(S_{3b}(n_6, n_{11}, n_8))$ (published) & $500$ & $3.2588326203533$ \\
$K_{4a}(S_{3b}(n_{11}, n_{11}, n_{11}))$ & $660$ & $3.2588323392139$ \\
\bottomrule
\end{tabular}
\end{center}

Only two distinct values stand above the published combination, so it is third by value; among
the $16\,384$ candidates individually it is ninth, because the top value is attained four times.
Those four are $S_{3b}$ and $S_{3h}$, each on $n_6$ and on its orientation reversal, and they do not
merely agree to the printed digits: the four seven-family vectors differ, in the $A$ and $D$
coordinates, yet $K_{4a}$ returns the same $288$-digit integer for all of them. The bound is a
property of the tree, not of the particular rule or orientation that realises it.

\paragraph{The gain is not monotone.}
Using the strongest block three times is \emph{worse} than the record. The improvement is a
property of the mixture, not of block quality, which is why it is invisible to an argument of the
form ``use the best available piece''.

\paragraph{Order matters, and by more than the margin being contested.}
On the published split, permuting the three arguments of $S_{3b}$ costs up to
$2.3 \times 10^{-6}$ -- more than half of the $4.6 \times 10^{-6}$ by which the heterogeneous
refinement improved on BPZ. A parameter with that much leverage is worth searching.

\paragraph{Iteration does not compound.}
Feeding the result back, $X \mapsto S_{3b}(n_6, X, X)$, gains at the first step and loses
thereafter: $3.2588343622$, then $3.2588110347$, then $3.2587875577$. The recursion has an
interior optimum, so depth alone is not the lever either. Likewise, applying a terminal code to
four \emph{distinct} representations was in every case we tried worse than four copies of the
best one, which is what one expects of a geometric mean.

\section{What does not help}

Several natural directions can be closed exhaustively rather than merely explored. We record them
because each cost real effort and each is cheap to state.

\paragraph{The base gadget admits exactly eight private pairs.}
For the $367$-set $I$ of Polak and Schrijver~\cite{PS} in $C_7^{\boxtimes 5}$, a private pair needs a vertex
$q \neq r$ with $N[q] \cap I = \{r\}$. Exactly eight members of $I$ own such a vertex, each owning
exactly one, and they are precisely the eight pairs in use. So $t = 8$ is a structural ceiling for
this set, not an artefact of search. At dimension ten the count is likewise exactly $5152$, which
is what the recursion already delivers.

\paragraph{The neutral part cannot exceed $322$ over the whole orbit.}
Write $U$ for the closed neighbourhood of the sixteen pair endpoints; then $o = |X| - |X \cap U|$,
and $U$ does not depend on the choice of transversals. One more neutral point, $o = 323$, would
suffice to beat the record by squaring alone. Applying every automorphism of $C_7^{\boxtimes 5}$
to $I$ while holding $X$ fixed -- all $5! \cdot 2^5 \cdot 7^5 = 64\,538\,880$ of them -- and
imposing the condition $X \cap N[P^H] \cap N[P^V] = \emptyset$, no image reaches $o = 323$.
A separate walk of four million size-preserving exchanges on $X$ at fixed $|X| = 367$ also never
exceeded $322$.

\paragraph{Nine private pairs exist but are unusable.}
Iterated local search produced $65$ distinct $367$-sets in $C_7^{\boxtimes 5}$. Their private-pair
counts run from $t = 5$ to $t = 10$: two, twenty-one, sixteen, nineteen, six and one respectively.
A ninth pair would be worth a great deal: at $s = 367$ any $o \ge 311$ beats the record, and the
nine-pair bases offer $315$ to $317$. They are nonetheless unusable, and the obstruction is
quantitative. The doubly covered region $D = N[P^H] \cap N[P^V]$, which the auxiliary set must
avoid, grows by almost exactly the same amount with each extra pair. Over all legal splits of
every base we hold, its size runs $956$ to $988$ for the nineteen eight-pair bases -- the
published $X$ dodges one of size $956$ -- then $1098$ to $1128$ for the six with nine pairs, and
$1240$ to $1268$ for the single base with ten. That is about $142$ further forbidden vertices per
pair, three times over, which makes the obstruction look structural rather than incidental, and
closes the ten-pair base for the same reason as the nine-pair ones. Sweeping the whole
automorphism group against a nine-pair base -- all $64\,538\,880$ images, so the entire orbit,
against thirty-two orbit representatives as auxiliary set -- no image satisfies the disjointness
condition. At $|X| = 366$ the pairs are
available but then $o \ge 338$ would be needed, which is out of reach; the question is binary and
it answers no.

\paragraph{The published twist is optimal.}
The heterogeneous step reads $q = |J \setminus N(X^0)|$ for a codebook $J$, and an automorphic
image of an independent set is again one, so the twist is free. Sweeping all $64\,538\,880$
automorphisms, the published choice maximises $q/|J|$ at dimensions $10$, $15$, $20$ and $25$ --
strictly, at every one of them.

\paragraph{The rule library is saturated where it matters.}
Admissibility is two local conditions, so one can ask whether the published rules can be extended.
No word can be added to the $B, N, A, D, H$ or $V$ component of any rule. The $O$ component is not
saturated -- $S_{3a}$ extends from four words to seven, $S_{3b}$ from none to four, five rules
from none to nine -- but this buys nothing, because no terminal code reads the label $O$, and no
component other than $O$ takes $O$ as input. Saturating $O$ and rerunning the published recursion
reproduces the published number to the last digit. The three terminal codes admit no extension at
all. Randomly generated admissible rules, including four-input rules, which the library does not
contain, are far worse than the published ones.

\paragraph{The library is not closed under the orientation involution.}
Exchanging $A \leftrightarrow D$ and $H \leftrightarrow V$ is an automorphism of the separation
relation, so the conjugate of an admissible rule is admissible. Only $S_{2a}$ and $S_{3a}$ are
self-conjugate: the conjugates of the other eight rules and of all three terminal codes are
admissible but absent. A search that tries both orientations of every input covers them
implicitly.

\section{Two traps that internal consistency does not catch}

Both cost us a day, and both produced numbers that looked like records.

\paragraph{The class rule is a theorem only for homogeneous products.}
Tracking which vertices lie in $N[X^0]$ under a product is needed to propagate $q$. For the
homogeneous product the rule
\[
  i(u,w) = (i_u \wedge i_w) \vee (j_u \wedge j_w), \qquad
  j(u,w) = (j_u \wedge i_w) \vee (i_u \wedge j_w)
\]
is a three-line theorem: $N[(u,w)] = N[u] \times N[w]$, a rectangle $A \times B$ is met exactly
when $u$ meets $N[A]$ and $w$ meets $N[B]$, and $X^0$ of the product is the union of two
rectangles. For the heterogeneous product the neutral part is built from a different rectangle,
whose left factor is the escaping part of the codebook, so the rule needs a different predicate
$E$. On the base gadget $E$ differs from $j$ on $5589$ of $16807$ vertices. The arithmetic of the
wrong rule is perfectly self-consistent, and the identity $s = o + h + v$ holds throughout.
Recomputing in exact integers does not help: exact arithmetic of a wrong formula stays wrong.

Moreover the predicate $E$ does not close in any bounded class space. Refining the partition of
$C_7^{\boxtimes 5}$ by the predicates it generates reaches a fixed point at $16\,776$ classes
among $16\,807$ vertices; the class degenerates to the vertex. This explains why $q_{10}$ and
$q_{15}$ appear in the literature as computed constants rather than as terms of a recursion.

\paragraph{Check against the upper bound.}
A search that allowed a fifteen-dimensional codebook to serve a ten-dimensional left gadget --
an independent set of $49\,495\,055$ vertices in a graph whose independence number is at most
$134\,753$ -- produced $4.11$, with every internal invariant satisfied. Only comparison with
$\vartheta(C_7) = 3.3176672$ exposed it. In a framework where all the arithmetic is sums of
products of non-negative integers, a wrong object does not announce itself; the upper bound is the
cheapest independent alarm available.

\section{Verification}

Everything was computed twice, by implementations sharing no code.

The first is in JavaScript. It parses the rule library from the Lean sources of the BPZ
repository at the commit Tandon names, checks that each rule carries the number of words the paper
states and that each satisfies both admissibility conditions, and implements \eqref{eq:card} in
arbitrary-precision integers.

The second is in Python. Its copies of $S_{3b}$ and $K_{4a}$ are typed from the printed tables of
Tandon's Appendix A, not from the machine-readable sources, and it implements the cardinality map
independently.

Both are checked against the published construction before being used on ours: each reproduces
Tandon's $n_{25}$ in all seven coordinates and his $257$-digit $M$ exactly, and hence his stated
bound. On the construction of Theorem~\ref{thm:main} they agree on a $288$-digit integer and on
\[
  M^{1/560} = 3.258834362237710794\ldots
\]
As a sanity check the value lies below $\vartheta(C_7)$, as any lower bound must.

Two further points of external agreement are worth recording, as they were obtained
independently of the paper and only afterwards compared with it. Computing the region split of
the dimension-ten gadget directly from its explicit vertex sets gives $h = 12\,236$,
$v = 16\,744$, which is what Tandon's Table 6 lists; and $|X_{10} \setminus N(X^0_{10})| =
134\,689 - 105\,709 = 28\,980$ likewise agrees.

\section*{Reproducibility and assistance}

The certificate, both verification scripts, the search of Section~4 and the census and
obstruction checks of Section~5 are in the accompanying repository, archived at
\texttt{doi:10.5281/zenodo.22727744}. The rule library is anchored
to commit
\texttt{aa21eeb12b75b041}\allowbreak\texttt{3d3fa9fb4208b5d0bf2c4d65} of the BPZ Lean repository, as in Tandon's
paper.

The search over recursion trees, the exhaustive verifications of Section~5 and the two
implementations were carried out with the assistance of a large language model. All results are
reproducible from the scripts and certificates supplied; the Python verification was written
independently of the search code and from the printed rule tables rather than from the
machine-readable sources, so that it does not share the assumptions of the code that produced the
result.

\begin{thebibliography}{9}
\bibitem{Shannon} C.~E. Shannon, The zero error capacity of a noisy channel,
  \emph{IRE Trans. Inform. Theory} \textbf{2} (1956), 8--19.
\bibitem{Lovasz} L.~Lov\'asz, On the Shannon capacity of a graph,
  \emph{IEEE Trans. Inform. Theory} \textbf{25} (1979), 1--7.
\bibitem{PS} S.~C. Polak and A.~Schrijver, New lower bound on the Shannon capacity of $C_7$ from
  circular graphs, \emph{Inform. Process. Lett.} \textbf{143} (2019), 37--40.
\bibitem{Itty} N.~Itty, C.~D. Rosin, C.~Carstensen and D.~Reichman, Improved lower bounds for the
  Shannon capacity of odd cycles, arXiv:2607.21517.
\bibitem{Gao} Y.~Gao, A recursive construction improving the lower bound on the Shannon capacity
  of $C_7$, arXiv:2607.27869.
\bibitem{BPZ1} P.~Buys, S.~Polak and J.~Zuiddam, Lean-verified lower bounds for the Shannon
  capacity of odd cycles, arXiv:2607.29681.
\bibitem{BPZ2} P.~Buys, S.~Polak and J.~Zuiddam, Lean-verified lower bounds for the Shannon
  capacity of odd cycles, updated version carrying the multi-gadget framework; Lean repository
  at commit \texttt{aa21eeb12b75b041}\allowbreak\texttt{3d3fa9fb4208b5d0bf2c4d65}.
\bibitem{Tandon} R.~Tandon, Strengthening recursive constructions for zero-error Shannon capacity,
  arXiv:2608.30273.
\end{thebibliography}

\appendix
\section{The certificate}

The cardinality vector $v = S_{3b}(n_6, n_{11}, n_{11})$ of a seven-family representation
in $C_7^{\boxtimes 140}$, in the order $(B,N,A,D,O,H,V)$. The zero in the $O$ coordinate
reflects $T^{(3b)}_O = \emptyset$, as in the published construction.

{\small\begin{verbatim}
F_B = 630253135055981401870713517244555479224707417188321589099630892251191137
F_N = 144695722184541768107720115813414797839355308000072965335422489005149364
F_A = 234597357441727401048883289845530873859970864445720627213945096783673569
F_D = 279954818401127974982259955129730445935176686603744292186306152889419948
F_O = 0
F_H = 43369967591254229190937078233276375615707839828850482473899658380862816
F_V = 43369967591254229190937078233276375615707839828850482473899658380862816
\end{verbatim}}

Applying $K_{4a}$ to four copies gives an independent set in $C_7^{\boxtimes 560}$ whose
cardinality has $288$ digits:

{\small\begin{verbatim}
20648129955823938354627223642277723610774954309518268857043903735102296965
39692406790071877020909818803770098048843510624763953396835623197404041433
74478908724579606022957152596859036287499885258150979930066394024808709850
012732304675637453738651340534438147137232028524828404364024830721
\end{verbatim}}

Hence
\[
  \Theta(C_7) \;\ge\; M^{1/560} \;=\; 3.258834362237710794\ldots,
\]
against $3.258832620353266309\ldots$ for the published construction. To thirty places the
value is $3.258834362237710794353262972358$.

\end{document}
